\section{Introduction}
\label{sec:introduction}

The rapid dissemination of scientific claims through social media poses a substantial challenge for fact-checking, as
posts frequently summarize, simplify, or sensationalize research findings while omitting an explicit reference to the
original study. As a result, the scientific source underlying a claim is rarely provided through a direct URL or a
complete bibliographic citation and is instead only hinted at through partial cues, such as the reported finding, the
publication venue, the research institution, or the authors involved, making verification not merely a matter of
assessing the plausibility of the claim itself, but first and foremost a matter of identifying the scientific
publication to which the post implicitly refers.

This problem is addressed in Task 1 of the CLEF 2026 CheckThat! Lab, \emph{Source Retrieval for Scientific Web Claims},
whose objective is to identify, for each social media post containing a scientific claim and an implicit reference to a
study, the corresponding publication within a fixed pool of 10,000 English-language candidate papers.

The task is particularly challenging because it requires bridging the marked gap between the language of social media
and that of scientific writing, as posts are typically short, informal, and often expressed through paraphrases,
slogans, hashtags, or incomplete statements, whereas scientific papers adopt a much more precise and technical register.
At the same time, the referenced source is not cited explicitly, meaning that the retrieval system cannot rely on direct
bibliographic signals and must instead infer the connection from partial semantic evidence; this challenge is further
compounded by the variety of forms in which relevant clues may appear, including mentions of a scientific finding, a
medical outcome, a research institution, or a publication venue, thus requiring the effective integration of lexical and
semantic evidence.

A further layer of complexity is introduced by the multilingual nature of the benchmark, since the official query sets
include posts in English, German, and French, while the candidate collection is shared across languages. Consequently,
an effective retrieval system must not only reconstruct the implicit link between a post and its source publication, but
must also do so in a setting where the evidence is expressed in multiple languages and must ultimately be mapped onto
the same English-language document pool.

Team Retrix addresses this task through a four-stage sequential pipeline (Figure~\ref{fig:pipeline}):
\begin{enumerate}
    \item {Queries originally written in French or German are first translated into English using
    \sloppy{utter-project/EuroLLM-9B-Instruct-2512}, so that all downstream components can operate over a unified input
    space independently of the source language.}

    \item {The resulting English queries are then enriched through a pseudo-relevance feedback strategy, where
    \sloppy{BAAI/bge-large-en-v1.5} is used to retrieve semantically related papers and extract salient expansion terms,
    thereby helping bridge the gap between the informal language of social media posts and the more specialized
    terminology of scientific publications.}

    \item {The expanded queries are subsequently processed by a hybrid retriever based on \sloppy{BAAI/bge-m3}, which
    combines dense, sparse, and ColBERT-style signals in order to generate a high-recall list of candidate papers from
    the shared collection.}

    \item {The retrieved candidates are reranked with a neural cross-encoder, with
    \sloppy{nvidia/llama-nemotron-rerank-1b-v2} providing our strongest single-stage configuration, while additional
    experiments with cascaded and fine-tuned rerankers were conducted to further improve the ordering of the top results.}
\end{enumerate}

Throughout the development of this pipeline, we also investigated the effect of fine-tuning reranking models on
domain-specific hard negatives derived from the training data. \\

%Todo : inserire il posizionamento
Evaluated on the CheckThat! 2026 Task 1 dataset our system ranked XXXXXXXXXXXXX on the official leaderboard attaining an
average Mean Reciprocal Rank at 5 (MRR@5) of XXXXXXXXXXXXXX \% on the development set (XXXXX place). \\

The remainder of this paper is organized as follows: Section~\ref{sec:related} reviews the research directions most
closely related to the task, Section~\ref{sec:methodology} presents our approach in detail, Section~\ref{sec:setup}
describes the experimental setup, and Section~\ref{sec:results} discusses the main results, while
Section~\ref{sec:conclusion} concludes the paper by summarizing our findings and outlining possible directions for
future work.
